\subsection{Inference when the direction is estimated}
\label{sec:estimated-direction-inference}

Theorem~\ref{thm:idprime} fixes the direction and asks what it identifies. In
practice the direction is not fixed; it is read off the data, as the leading
eigenvector of the residualized embedding covariance or as the concept-orthogonal
projection the erasure supplies, and an interval that treats it as known understates
the sampling variability. This section gives the limiting distribution of the
estimator that accounts for the estimated direction, under the spatial dependence the
panel carries.

Write $\hat v$ for the estimated direction, formed on the pooled sample of $N$
listings across the panel, and $\hat\theta_k=\hat v'\hat a_k/(\hat v'\hat\Sigma_k\hat
v)$ for the per-market estimate on $n_k$ listings, the sample analogue of
\eqref{eq:thetav-estimable} with $\hat a_k$, $\hat\Sigma_k$ the cross-fitted residual
moments in market $k$. Two facts shape the analysis. The estimand is not orthogonal to
the direction: the gradient
\[
\frac{\partial\theta}{\partial v}=\frac{a-2\,\theta(v)\,\Sigma v}{v'\Sigma v}
\]
does not vanish at the identified direction, so a Neyman-orthogonality argument cannot
make $\hat v$'s error disappear. What makes it negligible instead is a comparison of
rates. The direction is estimated on the pooled corpus, whose size $N$ exceeds any
single market's $n_k$ by construction, and its error enters $\hat\theta_k$ scaled by
that larger sample.

\begin{assumption}[Spectral gap]\label{as:ed-gap}
The pooled residualized embedding covariance has a bounded eigen-gap separating the
direction $v^\ast$ from the remainder of its spectrum.
\end{assumption}
\begin{assumption}[Rate domination and independence]\label{as:ed-rate}
$n_k/N\to0$ for each market, and $\hat v$ is formed out-of-fold, on a sample
independent of the observations used to form $\hat\theta_k$, in practice by leaving
market $k$ out of the pooled direction fit.
\end{assumption}
\begin{assumption}[Spatial weak dependence]\label{as:ed-spatial}
Within each market the listings occupy distinct locations separated by a minimum
distance, so the sampling region expands rather than fills in, and the
influence-function terms form an $\alpha$-mixing random field on geographic distance
whose mixing coefficients satisfy the summability conditions of
\citet{jenish2009central}; the cross-fitted nuisances satisfy the rates of
\citet{chernozhukov2018double}.
\end{assumption}

\begin{theorem}[Estimated-direction spatial DML]\label{thm:estimated-direction}
Under Assumptions~\ref{as:ed-gap}--\ref{as:ed-spatial}, for each market
\[
\sqrt{n_k}\,\bigl(\hat\theta_k-\theta_k(v^\ast)\bigr)\;\xrightarrow{d}\;\mathcal
N(0,V_k),
\]
where $\theta_k(v^\ast)$ is the direction-indexed estimand \eqref{eq:thetav-estimable}
and $V_k$ is the asymptotic variance of the Robinson influence function
$\psi_i=\mathbb{E}[\tilde T^2]^{-1}\,\tilde T_i(\tilde Y_i-\theta_k\tilde T_i)$ under
spatial dependence. The estimated direction contributes at $o_p(n_k^{-1/2})$, so $V_k$
is the fixed-direction DML variance, consistently estimated by a spatial
heteroskedasticity- and autocorrelation-consistent estimator of the
influence-function sum \citep{conley1999gmm,bester2011inference}.
\end{theorem}

\begin{proof}[Proof sketch]
The argument composes four steps. First, under the spectral gap of
Assumption~\ref{as:ed-gap} the estimated direction obeys a normal approximation for
bilinear forms of the sample spectral projector at rate $N^{-1/2}$
\citep{koltchinskii2016asymptotics}, so $\|\hat v-v^\ast\|=O_p(N^{-1/2})$; the
deterministic subspace bound of \citet{yu2015useful} certifies this rate but does not
supply a distribution, and is used only for the former. Second, a first-order
expansion of $\hat\theta_k$ in the direction gives
$\hat\theta_k=\theta_k(v^\ast)+(\partial\theta/\partial v)'(\hat v-v^\ast)+r_k$, and by
Assumption~\ref{as:ed-rate} the linear term is $O_p(N^{-1/2})=o_p(n_k^{-1/2})$, while
the out-of-fold construction makes $\hat v$ independent of the market-$k$ score, so no
covariance between the two enters the limiting variance. Third, with the direction held
at $v^\ast$, the cross-fitted Robinson estimator is asymptotically linear with
influence function $\psi_i$, the nuisance error entering at second order by Neyman
orthogonality \citep{chernozhukov2018double}. Fourth, the normalized sum of $\psi_i$
over the spatially dependent field is asymptotically normal by the central limit
theorem of \citet{jenish2009central}, with variance consistently estimated by the
spatial-HAC estimator. Composing the four steps, the direction's error is
asymptotically negligible and the limit is the fixed-direction spatial-DML normal law.
\end{proof}

Two remarks fix the scope. The negligibility of the direction's error is earned by the
rate comparison of Assumption~\ref{as:ed-rate}, not by orthogonality; were the
direction formed in sample, or were a single market as large as the pool, the
generated-regressor correction of \citet{pagan1984econometric} and
\citet{hahn2013asymptotic} would enter, and $V_k$ would carry an added term in the
asymptotic covariance of $\hat v$ that the spectral theory of
\citet{koltchinskii2016asymptotics} makes estimable. We avoid that term by the pooled,
out-of-fold construction, and say so rather than assume it away. And the result is what
\citet{guiveitch2023} lack: their central limit theorem for a text-treatment effect
handles a learned representation through the usual nuisance-rate condition, but their
treatment is not a direction, so it carries no term for an estimated direction entering
as a generated regressor. Theorem~\ref{thm:estimated-direction} supplies exactly that
term and shows it is second order, which is what licenses the per-market intervals the
panel reports once they are widened to the spatial-HAC scale.
Section~\ref{sec:identified-direction} carries this out on the identified direction,
where the out-of-fold construction reproduces the point estimates while the Conley
spatial-HAC intervals widen by a factor near two and the effect survives in nine of
the twelve markets.

\paragraph{What the paper identifies, and what it bounds.} The two theorems delimit
the claim the empirical sections make, and the paper holds to one line throughout.
Against the geography that motivates it, the effect is identified by design: the
identified direction removes the location the controls and the erasure can express,
the within-building rung holds location fixed by construction, and
Theorem~\ref{thm:idprime} names the direction on which the removal is exact while
Theorem~\ref{thm:estimated-direction} supplies the inference once that direction is
estimated. Against what the design cannot reach, unmeasured interior quality and the
nonlinear spatial leakage the linear erasure leaves behind, the effect is not
identified but bounded, by the robustness value of Section~\ref{sec:sensitivity-12}
and by the near-twin matching of Section~\ref{sec:within-building}. The estimate is
therefore a design-identified effect of the text direction against observed location,
carrying a sensitivity bound on the residual, and not a structural semantic parameter
free of every confound. Where the paper writes that the effect is identified it means
the first; where it writes that the effect is associational and sensitivity-bounded it
means the second; the two are the same claim read from its two sides.
